\chapter{Crossover, Weak Harnack, Harnack, and H\"older Regularity}\label{sec:crossover}
%% ========================================================================

%% --- Section 9 ---
\section{Crossover Estimate}

Let $d \in \N$, $d \geq 3$, and $E = \R^d$ throughout. Write $\Bz r$ for the
open Euclidean ball of radius $r$ centred at the origin. Given a normalized
elliptic coefficient $A$ on $\Bz 1$ with ellipticity constants $\lambda_A,
\Lambda_A$ (with $\lambda_A = 1$ after normalization, so $\Lambda_A \geq 1$),
write $\IsSupersolution{A}{u}$ for the assertion that $u \in W^{1,2}_{\loc}$
satisfies $-\dv(A\nabla u) \geq 0$ in the weak sense, and $\IsSolution{A}{u}$
for the corresponding equality.

The crossover estimate, due in this form to Moser, bridges positive and
negative powers of a positive supersolution. Throughout this section we work
with a fixed positive constant
\begin{equation*}
  \Cpoinc = \Cpoinc(d) > 0
\end{equation*}
denoting the dimensional Poincar\'e constant on the unit ball, a fixed
constant $\Mst > 1$ used in the smooth cutoff construction, and the
dimensional John--Nirenberg constant $\CJN(d) > 0$.

\begin{definition}
\label{def:c_crossover_bmo_scale}
\lean{DeGiorgi.c_crossover_bmo_scale, DeGiorgi.c_crossover'}
\leanok
\uses{def:C_poinc_val, def:C_JN, lem:smoothTransition_nnnorm_deriv_bounded}
We define the BMO scale and crossover exponent
\begin{align*}
  c_{\BMO}(d) &:= \bigl(\vol(\Bz 1)\bigr)^{-1/2}\,\Cpoinc(d) \cdot
    \tfrac12^{1 - d/2} \cdot 8\Mst \,\bigl(\vol(\Bz 1)\bigr)^{1/2}, \\
  c'(d)       &:= \frac{1}{2\,\CJN(d)\,c_{\BMO}(d) + 1}.
\end{align*}
\end{definition}

\begin{lemma}
\label{lem:c-crossover-bounds}
\lean{DeGiorgi.crossoverC'_pos, DeGiorgi.c_crossover'_lt_one}
\leanok
\uses{def:c_crossover_bmo_scale}
$0 < c'(d) < 1$ and $c'(d) \leq 1$.
\end{lemma}

\begin{proof}
\leanok
\uses{def:C_poinc_val, def:C_JN, lem:smoothTransition_nnnorm_deriv_bounded}

The constant $c_{\BMO}(d)$ is a product of strictly positive terms, so
$c_{\BMO}(d) > 0$. Since $\CJN(d) > 0$, the denominator
$2\,\CJN(d)\,c_{\BMO}(d) + 1$ is strictly greater than $1$, hence its
reciprocal lies in $(0,1)$.
\end{proof}

\subsection{Logarithmic gradient bound}

The first ingredient is the Caccioppoli-type inequality for $\log u$, obtained
by testing the supersolution against $\varphi^2/(u + \varepsilon)$.

\begin{theorem}
\label{thm:log-grad-eps}
\lean{DeGiorgi.log_gradient_bound_eps}
\leanok
\uses{def:MemW1pWitness, def:EllipticCoeff, def:IsSubsolution}
Let $\Omega \subseteq E$ be open, $A$ an elliptic coefficient on $\Omega$
with ellipticity constants $\lambda, \Lambda$, $u \in W^{1,2}_{\loc}(\Omega)$
positive on $\Omega$ with $\IsSupersolution A u$, and let
$\varphi \in C_c^\infty(\Omega)$. Then for every $\varepsilon > 0$,
\begin{equation*}
  \int_\Omega \frac{\varphi^2 \,|\nabla u|^2}{(u+\varepsilon)^2}\,dx
    \;\leq\; \frac{4\Lambda}{\lambda}\int_\Omega |\nabla\varphi|^2\,dx.
\end{equation*}
\end{theorem}

\begin{proof}
\leanok
\uses{def:HasWeakPartialDeriv, def:MemW1p, def:MemW01p, lem:witness_add, lem:MemW1pWitness_restrict, prop:witness_mul_smooth_bounded, lem:witness_smooth_cpt, prop:H01_vector_space, lem:memW01p_of_smooth_cpt, thm:memW01p_of_compactly_supported, def:smoothGradNorm, thm:chain-rule-open, lem:EllipticCoeff_lam_nonneg, lem:EllipticCoeff_ae_coercive_nonneg, def:bilinFormIntegrandOfCoeff, lem:aestronglyMeasurable_bilinFormIntegrandOfCoeff, def:IsSmoothTestOn, def:smoothGradField, thm:bilinForm_weighted_cauchySchwarz, lem:integrable_bounded_mul_bilinFormIntegrand}

Set $\psi := \varphi^2/(u+\varepsilon)$. By the smooth chain-rule witness
construction \leanname{build\_test\_function\_eps}, $\psi$ lies in
$H^1_0(\Omega)$, $\psi \geq 0$, and its weak gradient satisfies almost
everywhere
\begin{equation*}
  \nabla\psi = \frac{2\varphi\,\nabla\varphi}{u+\varepsilon}
    - \frac{\varphi^2\,\nabla u}{(u+\varepsilon)^2}.
\end{equation*}
Since $u$ is a supersolution and $\psi \geq 0$,
\begin{equation*}
  0 \;\leq\; \int_\Omega \langle A\nabla u, \nabla\psi\rangle\,dx
   \;=\; \int_\Omega \frac{2\varphi}{u+\varepsilon}\langle A\nabla u,
       \nabla\varphi\rangle\,dx
   \;-\; \int_\Omega \frac{\varphi^2}{(u+\varepsilon)^2}
       \langle A\nabla u, \nabla u\rangle\,dx.
\end{equation*}
Set
\begin{align*}
  Q &:= \int_\Omega \frac{\varphi^2}{(u+\varepsilon)^2}\langle A\nabla u,\nabla u\rangle\,dx, \\
  X &:= \int_\Omega \frac{\varphi}{u+\varepsilon}\langle A\nabla u,\nabla\varphi\rangle\,dx.
\end{align*}
The supersolution inequality reads $Q \leq 2X$. By Cauchy--Schwarz applied to
the bilinear form $\langle A\,\cdot,\cdot\rangle$,
\begin{equation*}
  X^2 \;\leq\; Q \cdot \int_\Omega \langle A\nabla\varphi,\nabla\varphi\rangle\,dx
    \;\leq\; Q \cdot \Lambda \int_\Omega |\nabla\varphi|^2\,dx.
\end{equation*}
Combining $Q \leq 2X$ with $X^2 \leq Q\,\Lambda\int|\nabla\varphi|^2$, the
elementary inequality $Q^2 \leq 4QP$ implies $Q \leq 4P$ where
$P := \Lambda \int_\Omega |\nabla\varphi|^2$ (see
\leanname{log\_gradient\_algebraic\_conclusion}). Hence
\begin{equation*}
  \lambda \int_\Omega \frac{\varphi^2|\nabla u|^2}{(u+\varepsilon)^2}\,dx
    \;\leq\; Q \;\leq\; 4\Lambda \int_\Omega |\nabla\varphi|^2\,dx,
\end{equation*}
and dividing by $\lambda > 0$ gives the claim.
\end{proof}

\begin{theorem}
\label{thm:log-grad}
\lean{DeGiorgi.log_gradient_bound_of_supersolution}
\leanok
\uses{def:MemW1pWitness, def:EllipticCoeff, def:IsSubsolution}
Under the same hypotheses, with $u > 0$ on $\Omega$,
\begin{equation*}
  \int_\Omega \frac{\varphi^2\,|\nabla u|^2}{u^2}\,dx
    \;\leq\; \frac{4\Lambda}{\lambda}\int_\Omega |\nabla\varphi|^2\,dx.
\end{equation*}
\end{theorem}

\begin{proof}
\leanok
\uses{thm:log-grad-eps}

Apply Theorem~\ref{thm:log-grad-eps} with $\varepsilon_n = 1/(n+1)$. The
integrand
\begin{equation*}
  g_n(x) := \varphi(x)^2|\nabla u(x)|^2/(u(x)+\varepsilon_n)^2
\end{equation*}
increases monotonically as $\varepsilon_n \downarrow 0$ to
$g_0(x) := \varphi(x)^2|\nabla u(x)|^2/u(x)^2$ at every point of $\Omega$
since $u > 0$ there. By Fatou's lemma,
\begin{equation*}
  \int_\Omega g_0\,dx \;\leq\; \liminf_{n\to\infty}\int_\Omega g_n\,dx
    \;\leq\; \frac{4\Lambda}{\lambda}\int_\Omega |\nabla\varphi|^2\,dx.
\end{equation*}
\end{proof}

\subsection{BMO control of \texorpdfstring{$\log u$}{log u}}

For $\varepsilon > 0$ define the regularized log
$v_\varepsilon(x) := -\log(u(x)+\varepsilon) + \log\varepsilon$ and its smooth
chain-rule witness $\leanname{regularizedLogFun}$. Combining the log-gradient
bound with the Poincar\'e inequality on $\Bz{1/2}$ gives the BMO estimate.

\begin{theorem}
\label{thm:bmo-half}
\lean{DeGiorgi.regularizedLog_unit_halfBall_meanOscillation}
\leanok
\uses{def:ellipticityRatio, def:IsSubsolution}
Let $A$ be normalized elliptic on $\Bz 1$, $u > 0$ a supersolution. Then
for $v := v_\varepsilon$,
\begin{equation*}
  \fint_{\Bz{1/2}} \Bigl|v(x) - \fint_{\Bz{1/2}} v\Bigr|\,dx
    \;\leq\; c_{\BMO}(d)\,\Lambda_A^{1/2}.
\end{equation*}
\end{theorem}

\begin{proof}
\leanok
\uses{def:HasWeakPartialDeriv, def:MemW1p, def:MemW1pWitness, lem:weak_grad_norm_memLp, lem:zero_outside_tsupport, def:C_poinc_val, thm:poincare_unitBall_W1p_public, lem:eLpNorm_const_average_le, def:CampanatoBall, def:HasCampanatoBound, thm:chain-rule-open, lem:smoothTransition_nnnorm_deriv_bounded, thm:Cutoff_exists, lem:rescale-witness, lem:eLpNorm_rescale_to_unitBall, def:EllipticCoeff, def:c_crossover_bmo_scale, thm:log-grad-eps}

Pick a cutoff $\eta \in C_c^\infty(\Bz{3/4})$ with $\eta = 1$ on $\Bz{1/2}$,
$0 \leq \eta \leq 1$, and $|\nabla \eta| \leq 8\Mst$ (existence given by
\leanname{Cutoff.exists}). Theorem~\ref{thm:log-grad-eps} (with
$\varphi = \eta$, $\lambda_A = 1$) yields
\begin{equation*}
  \int_{\Bz{1/2}} |\nabla v_\varepsilon|^2\,dx
    \;\leq\; \int_{\Bz 1} \eta^2 \frac{|\nabla u|^2}{(u+\varepsilon)^2}\,dx
    \;\leq\; 4\Lambda_A \int_{\Bz 1}|\nabla \eta|^2\,dx
    \;\leq\; 4\Lambda_A \cdot (8\Mst)^2 \vol(\Bz 1).
\end{equation*}
Hence $\|\nabla v_\varepsilon\|_{L^2(\Bz{1/2})} \leq 16\Mst\,\Lambda_A^{1/2}\,\vol(\Bz 1)^{1/2}$.
By the Poincar\'e--Wirtinger inequality on $\Bz{1/2}$ with constant
$\Cpoinc(d) \cdot (1/2)^{1-d/2}$ (the dimensional rescaling factor),
\begin{multline*}
  \fint_{\Bz{1/2}}\Bigl|v_\varepsilon - \fint_{\Bz{1/2}} v_\varepsilon\Bigr|\,dx
    \;\leq\; \vol(\Bz{1/2})^{-1/2} \Cpoinc(d) (1/2)^{1-d/2}
        \|\nabla v_\varepsilon\|_{L^2(\Bz{1/2})} \\
    \;\leq\; \vol(\Bz 1)^{-1/2}\Cpoinc(d) (1/2)^{1-d/2}\cdot
        8\Mst\,\Lambda_A^{1/2}\,\vol(\Bz 1)^{1/2}
    \;=\; c_{\BMO}(d)\,\Lambda_A^{1/2}.
\end{multline*}
\end{proof}

The same statement holds on every ball $B(x_0, R/2)$ with $\overline{B(x_0,R)}
\subset \Bz 1$, by affine rescaling $\leanname{rescaleToUnitBall}$
(see \leanname{regularizedLog\_ball\_meanOscillation}).

\subsection{John--Nirenberg exponential integrability}

\begin{definition}
\label{def:JN-constants}
\lean{DeGiorgi.crossoverJNTheta, DeGiorgi.crossoverJNKloc, DeGiorgi.crossoverJNKhalf}
\leanok
\uses{def:C_JN, def:c_crossover_bmo_scale}
Set
\begin{align*}
  \theta(d)  &:= 24 \cdot 2^{d+1} \cdot \bigl(c'(d)\,c_{\BMO}(d)\bigr), \\
  K_{\loc}(d)&:= \exp(\theta(d)) \cdot (1 + \CJN(d)),\\
  K_{1/2}(d) &:= \lceil 97^d \rceil \cdot K_{\loc}(d).
\end{align*}
\end{definition}

\begin{theorem}
\label{thm:exp-half}
\lean{DeGiorgi.regularizedLog_halfBall_exp_average_to_origin_le}
\leanok
\uses{def:ellipticityRatio, def:IsSubsolution}
With $p := c'(d)/\Lambda_A^{1/2}$, $v := v_\varepsilon$, and
$a := \fint_{\Bz{1/48}} v\,dz$,
\begin{equation*}
  \fint_{\Bz{1/2}} \exp(p\,|v(z) - a|)\,dz \;\leq\; K_{1/2}(d).
\end{equation*}
\end{theorem}

\begin{proof}
\leanok
\uses{def:HasWeakPartialDeriv, def:MemW1p, def:MemW1pWitness, def:C_poinc_val, def:C_JN, thm:john_nirenberg_local, thm:chain-rule-open, lem:smoothTransition_nnnorm_deriv_bounded, lem:MemW1pWitness_sub_const, thm:finite-cover-card, def:EllipticCoeff, def:rescaleToUnitBall, def:rescaleCoeffToUnitBall, thm:rescaleToUnitBall_isSubsolution, def:EllipticCoeff_restrict, def:c_crossover_bmo_scale, lem:c-crossover-bounds, thm:bmo-half}

Apply the local John--Nirenberg lemma at scale $\rho = 1/48$ on a finite cover
of $\Bz{1/2}$ by balls $B(x_i, 1/48)$ with $\#\{x_i\} \leq \lceil 97^d\rceil$.
On each cover ball, Theorem~\ref{thm:bmo-half} (rescaled to $B(x_i, R)$ via
\leanname{regularizedLog\_ball\_meanOscillation}) gives
$\|v - \langle v\rangle\|_{\BMO} \leq c_{\BMO}(d)\,\Lambda_A^{1/2}$. The
classical John--Nirenberg theorem $\leanname{C\_JN}$ then yields
\begin{equation*}
  \fint_{B(x_i,1/48)} \exp\bigl(p\,|v - \langle v\rangle_{B(x_i,1/48)}|\bigr)\,dz
    \;\leq\; 1 + \CJN(d).
\end{equation*}
The shift between local and global averages is controlled by $\theta(d) =
24 \cdot 2^{d+1}\,c'(d)\,c_{\BMO}(d)$ via the standard telescoping
$\leanname{regularizedLog\_smallBallAverage\_to\_origin\_le}$. Multiplying
each local exponential bound by $\exp(p\,\theta(d)) = \exp(\theta(d))$ (since
$p\cdot \delta = \theta(d)$ by direct computation, using
$\delta := 24\cdot 2^{d+1}\,(c_{\BMO}(d)\Lambda_A^{1/2})$ and
$p \cdot \Lambda_A^{1/2} = c'(d)$), and summing over the $\lceil 97^d\rceil$
cover balls gives the asserted bound by $K_{1/2}(d)$.
\end{proof}

\subsection{Crossover product bound}

\begin{theorem}
\label{thm:crossover-exists}
\lean{DeGiorgi.crossover_product_bound_exists}
\leanok
\uses{def:ellipticityRatio, def:IsSubsolution, def:c_crossover_bmo_scale}
There exists a constant $C(d) \geq 1$ such that for every normalized elliptic
$A$ on $\Bz 1$, every positive supersolution $u$, and $p_0 := c'(d)/\Lambda_A^{1/2}$,
\begin{equation*}
  \Bigl(\fint_{\Bz{1/2}} u^{p_0}\,dx\Bigr)\Bigl(\fint_{\Bz{1/2}} u^{-p_0}\,dx\Bigr)
    \;\leq\; C(d).
\end{equation*}
\end{theorem}

\begin{proof}
\leanok
\uses{def:HasWeakPartialDeriv, def:MemW1p, def:MemW1pWitness, def:C_poinc_val, def:C_JN, thm:chain-rule-open, lem:smoothTransition_nnnorm_deriv_bounded, def:EllipticCoeff, lem:c-crossover-bounds, def:JN-constants, thm:exp-half}

By a regularization step $u \mapsto u + \varepsilon$, it suffices to bound
\begin{equation*}
  J_\varepsilon := \Bigl(\fint_{\Bz{1/2}}(u+\varepsilon)^{p_0}\Bigr)
        \Bigl(\fint_{\Bz{1/2}}(u+\varepsilon)^{-p_0}\Bigr)
\end{equation*}
uniformly in $\varepsilon > 0$ by $K_{1/2}(d)^2$. With $v_\varepsilon =
-\log(u+\varepsilon) + \log\varepsilon$, $a := \fint_{\Bz{1/48}} v_\varepsilon$,
$F_\pm(x) := \exp(\pm p_0(v_\varepsilon - a))$ and
$F_{\mathrm{abs}} := \exp(p_0|v_\varepsilon - a|)$, we have $F_\pm \leq
F_{\mathrm{abs}}$ pointwise. Theorem~\ref{thm:exp-half} together with
$\leanname{setAverage\_mono\_of\_nonneg\_local}$ yields
\begin{equation*}
  \fint_{\Bz{1/2}} F_+\,dx \leq K_{1/2}(d), \qquad
  \fint_{\Bz{1/2}} F_-\,dx \leq K_{1/2}(d).
\end{equation*}
Now $(u+\varepsilon)^{p_0} = \exp(p_0(\log\varepsilon - a))\,F_+$ and
$(u+\varepsilon)^{-p_0} = \exp(p_0(a - \log\varepsilon))\,F_-$ pointwise on
$\Bz{1/2}$. The two scalar prefactors multiply to $1$ (cancellation
$\leanname{hcancel}$), so
\begin{equation*}
  J_\varepsilon \;=\; \Bigl(\fint F_+\Bigr)\Bigl(\fint F_-\Bigr)
    \;\leq\; K_{1/2}(d)^2.
\end{equation*}
Letting $\varepsilon \downarrow 0$ along $\varepsilon_n = 1/(n+1)$, the
shifted integrals converge to the original ones by dominated/monotone
convergence (the integrability of $|u|^{p_0}$ is supplied by the
$W^{1,2}$ estimate using $p_0 < 1 \leq 2$, see
\leanname{hpow\_int}). Choose $C(d) := \max\{1, K_{1/2}(d)^2\}$.
\end{proof}

\begin{definition}
\label{def:Ccrossover}
\lean{DeGiorgi.C_crossover'}
\leanok
\uses{thm:crossover-exists}
$C_{\mathrm{cross}}(d) := C(d)$ from Theorem~\ref{thm:crossover-exists}, so
$C_{\mathrm{cross}}(d) \geq 1$.
\end{definition}

\begin{theorem}
\label{thm:crossover-estimate}
\lean{DeGiorgi.crossover_estimate}
\leanok
\uses{def:ellipticityRatio, def:IsSubsolution, def:c_crossover_bmo_scale, def:Ccrossover}
Let $d \geq 3$, $A$ normalized elliptic on $\Bz 1$, and $u > 0$ a
supersolution. Then with $c := c'(d)/\Lambda_A^{1/2}$,
\begin{equation*}
  \Bigl(\fint_{\Bz{1/2}} u^c\,dx\Bigr)\Bigl(\fint_{\Bz{1/2}} u^{-c}\,dx\Bigr)
    \;\leq\; C_{\mathrm{cross}}(d).
\end{equation*}
\end{theorem}

\begin{proof}
\leanok
\uses{thm:crossover-exists}

Set $v(x) := -\log u(x)$, $v_{\mathrm{avg}} := \fint_{\Bz{1/2}} v\,dx$,
$w(x) := \exp(c(v_{\mathrm{avg}} - v(x)))$, and
$w^{-1}(x) := \exp(c(v(x) - v_{\mathrm{avg}}))$, so $w \cdot w^{-1} \equiv 1$.
On $\Bz{1/2}$ where $u > 0$, $u^c = e^{-c v_{\mathrm{avg}}}\,w$ and
$u^{-c} = e^{c v_{\mathrm{avg}}}\,w^{-1}$, hence
\begin{equation*}
  \Bigl(\fint u^c\Bigr)\Bigl(\fint u^{-c}\Bigr)
    \;=\; e^{-cv_{\mathrm{avg}}}e^{cv_{\mathrm{avg}}}
        \Bigl(\fint w\Bigr)\Bigl(\fint w^{-1}\Bigr)
    \;=\; \Bigl(\fint w\Bigr)\Bigl(\fint w^{-1}\Bigr).
\end{equation*}
Apply Theorem~\ref{thm:crossover-exists}.
\end{proof}

\begin{corollary}
\label{cor:crossover-unaveraged}
\lean{DeGiorgi.crossover_estimate_unaveraged}
\leanok
\uses{def:ellipticityRatio, def:IsSubsolution, def:c_crossover_bmo_scale, def:Ccrossover}
Multiplying through by $\vol(\Bz{1/2})^2 > 0$,
\begin{equation*}
  \Bigl(\int_{\Bz{1/2}} u^c\Bigr)\Bigl(\int_{\Bz{1/2}} u^{-c}\Bigr)
    \;\leq\; C_{\mathrm{cross}}(d)\,\vol(\Bz{1/2})^2.
\end{equation*}
\end{corollary}

\begin{proof}
\leanok
\uses{thm:crossover-estimate}
See the Lean formalization.
\end{proof}

%% --- Section 10 ---
\section{Weak Harnack Inequality}

\begin{definition}
\label{def:p0}
\lean{DeGiorgi.weakHarnackP0, DeGiorgi.weak_harnack_decay_exp}
\leanok
\uses{def:ellipticityRatio, def:moserChi, def:c_crossover_bmo_scale}
For a normalized elliptic $A$ on $\Bz 1$, set
\begin{equation*}
  p_0 := p_0(A) := \frac{c'(d)}{\Lambda_A^{1/2}}, \qquad
  \beta(d) := \frac{d\,\chi(d)}{c'(d)},
\end{equation*}
where $\chi(d) := d/(d-2)$ is the Sobolev exponent (\leanname{moserChi}).
\end{definition}

\begin{lemma}
\label{lem:p_pos}
\lean{DeGiorgi.p₀_pos, DeGiorgi.p₀_lt_one, DeGiorgi.Λ_mul_p₀_sq}
\leanok
\uses{def:ellipticityRatio, def:c_crossover_bmo_scale, def:p0}
$0 < p_0(A) < 1$, $p_0(A) \leq c'(d)$, and $\Lambda_A\,p_0(A)^2 = c'(d)^2$.
\end{lemma}

\begin{proof}
\leanok
\uses{def:C_poinc_val, def:C_JN, lem:smoothTransition_nnnorm_deriv_bounded, lem:c-crossover-bounds}

Since $\Lambda_A \geq 1$, $\Lambda_A^{1/2} \geq 1$, so
$0 < p_0 = c'(d)/\Lambda_A^{1/2} \leq c'(d) < 1$. The identity follows from
$(c'(d)/\Lambda_A^{1/2})^2 \cdot \Lambda_A = c'(d)^2$.
\end{proof}

The proof of the weak Harnack inequality chains three pieces at half scale.
We rescale $v(z) := u(z/2)$ from $\Bz{1/2}$ to $\Bz 1$ via the rescaled
coefficient $A' := \leanname{rescaledCoeff}\,A$ which preserves $\Lambda$.
The three steps are:

\begin{itemize}
\item \emph{Forward Moser at level $p_0$}
  (\leanname{weak\_harnack\_stage\_one\_forward\_ball}):
\begin{equation*}
\Bigl(\int_{\Bz{1/4}} u^{q\,d/(d-2)}\Bigr)^{p_0(d-2)/(qd)}
\leq C_{\mathrm{wH},0\to}(d)\,(\Lambda_A p_0^2/(1-q)^2 + 1)^{d\chi/2}
\int_{\Bz{1/2}} u^{p_0}.
\end{equation*}

\item \emph{Crossover} (Corollary~\ref{cor:crossover-unaveraged}) at the level
$p_0$:
\begin{equation*}
\Bigl(\int_{\Bz{1/2}} u^{p_0}\Bigr)\Bigl(\int_{\Bz{1/2}} u^{-p_0}\Bigr)
\leq C_{\mathrm{cross}}(d)\,\vol(\Bz{1/2})^2.
\end{equation*}

\item \emph{Inverse Moser}
  (\leanname{weak\_harnack\_stage\_one\_inverse\_ball}):
\begin{equation*}
(\Lambda_A p_0^2 + 1)^{-d/2}\Bigl(\int_{\Bz{1/2}} u^{-p_0}\Bigr)^{-1}
\leq C_{\mathrm{wH},0}(d)\,\bigl(\essinf_{\Bz{1/4}} u\bigr)^{p_0}.
\end{equation*}
\end{itemize}

Multiplying these three estimates yields the chain.

\begin{theorem}
\label{thm:wH-chain}
\lean{DeGiorgi.weak_harnack_chain}
\leanok
\uses{def:ellipticityRatio, def:IsSubsolution}
For $d \geq 3$, $A$ normalized on $\Bz 1$, $u > 0$ a supersolution, and
$0 < q < 1$:
\begin{equation*}
\Bigl(\int_{\Bz{1/4}} u^{q d/(d-2)}\,dx\Bigr)^{p_0(d-2)/(qd)}
\leq \frac{C_{\mathrm{chain}}(d)}{(1-q)^{d\chi(d)}}
\Bigl(\essinf_{\Bz{1/4}} u\Bigr)^{p_0}.
\end{equation*}
\end{theorem}

\begin{proof}
\leanok
\uses{def:MemW1p, def:MemW1pWitness, def:C_poinc_val, def:C_JN, lem:smoothTransition_nnnorm_deriv_bounded, def:rescaleCoeffToUnitBall, thm:rescaleToUnitBall_isSubsolution, def:EllipticCoeff_restrict, lem:essSup_rescale_halfBall, def:C_MoserAnchor, def:moserChi, lem:moserDecayRatio_nonneg, def:C_Moser, lem:C_MoserAnchor_le_C_Moser, thm:preMoser-fwd, thm:wH-stage1-inv, def:C_weakHarnack0Forward, thm:wH-stage1-fwd, def:c_crossover_bmo_scale, lem:c-crossover-bounds, thm:crossover-exists, def:Ccrossover, cor:crossover-unaveraged, def:p0, lem:p_pos, lem:meas-bd}

We split into the case $p_0 < q$ (the natural case) and the case $q \leq p_0$
(handled by interpolating with $q' := (p_0+1)/2$). Assume $p_0 < q$ first.

By regularization $u \mapsto u + \delta$ (which preserves the supersolution
property by \leanname{isSupersolution\_add\_const\_unitBall}), and using the
identity $\Lambda_A p_0^2 = c'(d)^2$, the product of the forward Moser
estimate, the crossover bound and the inverse Moser estimate gives
\begin{multline*}
  \Bigl(\int_{\Bz{1/4}}(u+\delta)^{q d/(d-2)}\Bigr)^{p_0(d-2)/(q d)}
  \times \bigl(c'(d)^2 + 1\bigr)^{d/2}
  \\[-2pt]
  \leq C_{\mathrm{wH},0\to}\,C_{\mathrm{wH},0}\,C_{\mathrm{cross}}\,
       \vol(\Bz{1/2})^2\,(c'(d)^2/(1-q)^2 + 1)^{d\chi/2}
       \bigl(\essinf_{\Bz{1/4}}(u+\delta)\bigr)^{p_0}.
\end{multline*}
By \leanname{weak\_harnack\_forward\_factor\_bound} the right factor is
absorbed into $(1-q)^{-d\chi}$ times a dimensional constant, so collecting
terms produces a constant $C_{\mathrm{chainMain}}(d)$ with
\begin{equation*}
  \Bigl(\int_{\Bz{1/4}} (u+\delta)^{qd/(d-2)}\Bigr)^{p_0(d-2)/(qd)}
  \;\leq\; \frac{C_{\mathrm{chainMain}}(d)}{(1-q)^{d\chi(d)}}
        \bigl(\essinf_{\Bz{1/4}}(u+\delta)\bigr)^{p_0}.
\end{equation*}
Letting $\delta = 1/(n+1) \downarrow 0$ and using the dominated/monotone
convergence theorems on each side concludes the case $p_0 < q$.

For $q \leq p_0$ pick $q' := (p_0+1)/2 \in (p_0,1)$ and apply the previous
case with $q'$. Since $q d/(d-2) \leq q' d/(d-2)$, monotonicity of $L^p$ norms
on $\Bz{1/4}$ (\leanname{quarterBall\_lpNorm\_mono}) raises the resulting
$L^{q' d/(d-2)}$ norm to dominate the $L^{q d/(d-2)}$ norm, and a uniform
bound $C_{\mathrm{chainMain}}/(1-q')^{d\chi} \leq C_{\mathrm{chain}}/(1-q)^{d\chi}$
\leanname{(hmain\_to\_public)} then yields the claim.
\end{proof}

\begin{definition}
\label{def:CwH}
\lean{DeGiorgi.C_weakHarnack}
\leanok
\uses{def:moserChi, def:C_Moser, def:C_weakHarnack0Forward, def:c_crossover_bmo_scale, def:Ccrossover}
$
  C_{wH}(d) \;:=\; C_{\mathrm{chain}}(d)^{1/c'(d)}.
$
By Lemma~\ref{lem:c-crossover-bounds} and $C_{\mathrm{chain}} \geq 1$, we have
$C_{wH}(d) \geq 1$ (\leanname{one\_le\_C\_weakHarnack}).
\end{definition}

\begin{theorem}
\label{thm:wH}
\lean{DeGiorgi.weak_harnack}
\leanok
\uses{def:ellipticityRatio, def:IsSubsolution, def:p0, def:CwH}
For $d \geq 3$, $A$ normalized elliptic on $\Bz 1$, $u > 0$ a supersolution,
and $0 < q < 1$:
\begin{equation*}
  \Bigl(\int_{\Bz{1/4}} u^{q d/(d-2)}\,dx\Bigr)^{(d-2)/(qd)}
    \leq \Bigl(\frac{C_{wH}(d)}{(1-q)^{\beta(d)}}\Bigr)^{\Lambda_A^{1/2}}
        \essinf_{\Bz{1/4}} u.
\end{equation*}
\end{theorem}

\begin{proof}
\leanok
\uses{def:C_poinc_val, def:C_JN, lem:smoothTransition_nnnorm_deriv_bounded, def:moserChi, lem:moserDecayRatio_nonneg, def:C_Moser, lem:C_MoserAnchor_le_C_Moser, def:C_weakHarnack0Forward, def:c_crossover_bmo_scale, lem:c-crossover-bounds, thm:crossover-exists, def:Ccrossover, lem:p_pos, thm:wH-chain, lem:meas-bd}

By Theorem~\ref{thm:wH-chain},
$I^\alpha \leq C_{\mathrm{chain}}(d)/(1-q)^{d\chi}\cdot (\essinf u)^{p_0}$
where $I := \int_{\Bz{1/4}} u^{qd/(d-2)}$ and $\alpha := p_0(d-2)/(qd)$.
Raising both sides to the $1/p_0$ power (\leanname{rpow\_le\_rpow\_of\_exponent}):
\begin{equation*}
  I^{(d-2)/(qd)}
   \;\leq\; \Bigl(\frac{C_{\mathrm{chain}}(d)}{(1-q)^{d\chi(d)}}\Bigr)^{1/p_0}
       \essinf_{\Bz{1/4}} u.
\end{equation*}
Now use $1/p_0 = \Lambda_A^{1/2}/c'(d)$ and the laws of real powers:
\begin{align*}
  \Bigl(\frac{C_{\mathrm{chain}}(d)}{(1-q)^{d\chi(d)}}\Bigr)^{1/p_0}
   &= \Bigl(\bigl(C_{\mathrm{chain}}(d)/(1-q)^{d\chi(d)}\bigr)^{1/c'(d)}\Bigr)^{\Lambda_A^{1/2}}\\
   &= \Bigl(C_{\mathrm{chain}}(d)^{1/c'(d)}\big/(1-q)^{d\chi(d)/c'(d)}\Bigr)^{\Lambda_A^{1/2}}\\
   &= \Bigl(C_{wH}(d)\big/(1-q)^{\beta(d)}\Bigr)^{\Lambda_A^{1/2}}.
\end{align*}
\end{proof}

\begin{theorem}
\label{thm:wH-on-ball}
\lean{DeGiorgi.weak_harnack_on_ball}
\leanok
\uses{def:ellipticityRatio, def:IsSubsolution, def:p0, def:CwH}
On a general ball $\Ball R{x_0}$ with $R > 0$, the weak Harnack inequality
takes the form
\begin{multline*}
  \Bigl(R^{-d}\int_{\Ball{R/4}{x_0}} u^{qd/(d-2)}\,dx\Bigr)^{(d-2)/(qd)}\\
   \leq \Bigl(\frac{C_{wH}(d)}{(1-q)^{\beta(d)}}\Bigr)^{\Lambda_A^{1/2}}
        \essinf_{z \in \Bz{1/4}}\, u(x_0 + R z).
\end{multline*}
\end{theorem}

\begin{proof}
\leanok
\uses{def:rescaleToUnitBall, def:rescaleCoeffToUnitBall, thm:rescaleToUnitBall_isSubsolution, thm:wH}

Apply Theorem~\ref{thm:wH} to the rescaled function
$u_R(z) := u(x_0 + R z)$ on $\Bz 1$ with the rescaled coefficients
$A_R := \leanname{rescaleNormalizedCoeffToUnitBall}\,A$, which preserves
$\Lambda$ (so $p_0(A_R) = p_0(A)$). The change of variables
$\leanname{crossover\_integral\_comp\_affine\_ball}$ yields the
$R^{-d}$ Jacobian factor on the left, while the essential infimum of $u_R$
on $\Bz{1/4}$ equals that of $u$ on $\Ball{R/4}{x_0}$ pulled back through the
affine map.
\end{proof}

\subsection{Support constants}

\begin{definition}
\label{def:local-bases}
\lean{DeGiorgi.localMoserBase, DeGiorgi.localWeakHarnackBase}
\leanok
\uses{def:C_Moser, def:p0, def:CwH}
\begin{align*}
  M_{\mathrm{loc}}(d) &:= C_{\mathrm{Moser}}(d)\,((d-1))^d\,4^d, \\
  W_{\mathrm{loc}}(d) &:= C_{wH,0}(d)\, d^{\beta(d)}.
\end{align*}
\end{definition}

\begin{lemma}
\label{lem:one_le_localMoserBase}
\lean{DeGiorgi.one_le_localMoserBase, DeGiorgi.one_le_localWeakHarnackBase}
\leanok
\uses{def:local-bases}
For $d \geq 3$, $1 \leq M_{\mathrm{loc}}(d)$ and $1 \leq W_{\mathrm{loc}}(d)$.
Both are strictly positive.
\end{lemma}

\begin{proof}
\leanok
\uses{def:C_poinc_val, def:C_JN, lem:smoothTransition_nnnorm_deriv_bounded, def:moserChi, lem:moserDecayRatio_nonneg, def:C_Moser, lem:C_MoserAnchor_le_C_Moser, def:C_weakHarnack0Forward, def:c_crossover_bmo_scale, lem:c-crossover-bounds, thm:crossover-exists, def:Ccrossover, def:p0, def:CwH}

Each is a product of factors $\geq 1$: $C_{\mathrm{Moser}}(d) \geq 1$,
$(d-1)^d \geq 1$ (since $d-1 \geq 2$), $4^d \geq 1$, $C_{wH,0} \geq 1$,
$d^{\beta(d)} \geq 1$ (since $d \geq 1$ and $\beta(d) \geq 0$ by
\leanname{weak\_harnack\_decay\_exp\_nonneg}).
\end{proof}

\subsection{Measure-bound utilities}

The Harnack and Holder layers also use a small toolkit of $\esssup/\essinf$
arithmetic on restricted ball measures (\leanname{Support/MeasureBounds.lean}).

\begin{lemma}
\label{lem:meas-bd}
\lean{DeGiorgi.le_essInf_real_of_ae_le, DeGiorgi.essSup_le_of_ae_bdd, DeGiorgi.essInf_add_const_of_bdd, DeGiorgi.essSup_neg_of_bdd}
\leanok
Let $\mu$ be a non-zero (Borel) measure on $E$ and $u : E \to \R$.
\begin{enumerate}
\item If $u \geq c$ $\mu$-a.e., then $c \leq \essinf_\mu u$.
\item If $K \leq u \leq C$ $\mu$-a.e., then $\essinf_\mu u, \esssup_\mu u
  \in [K,C]$.
\item If $u$ is bounded $\mu$-a.e., then
  $\essinf_\mu(u+c) = \essinf_\mu u + c$ and the symmetric statement for
  $\esssup$.
\item $\esssup_\mu(-u) = -\essinf_\mu u$ when $u$ is bounded.
\item For a measurable embedding $T : E\to E$ and measurable $g$,
  $\essinf_\mu(g\circ T) = \essinf_{T_*\mu} g$ \leanname{(essInf\_comp\_measurableEmbedding\_eq)}.
\item For $c \in \ENNReal\setminus\{0\}$, $\essinf_{c\cdot \mu}f =
  \essinf_\mu f$.
\end{enumerate}
Furthermore, $\vol(\Ball r c)\neq 0$ for $r > 0$, so the restriction
$\vol|_{\Ball r c}$ is non-zero \leanname{(restrict\_ball\_ne\_zero)}.
\end{lemma}

\begin{proof}
\leanok

Each statement is a direct unfolding of $\essinf$/$\esssup$ as a $\sSup$/$\sInf$
over levels at which the corresponding sublevel set has positive/full measure.
For instance, in (1): assume $\mu \neq 0$ and $\mu(\{u < c\}) = 0$. Then
$\{u < a\}$ has measure zero for every $a < c$, so $a$ is dominated by every
upper bound used in the definition of $\essinf_\mu u$, hence $c \leq
\essinf_\mu u$. Statements (3),(4) follow from the corresponding
$\liminf/\limsup$ arithmetic on the filter $\ae \mu$ together with the fact
that the functions are eventually bounded. Statements (5),(6) follow from
the change-of-variables and scaling formulas for the underlying
$\sSup/\sInf$ over level sets.
\end{proof}

%% --- Section 11 ---
\section{Harnack Inequality and Holder Regularity}

\begin{definition}
\label{def:Charnack}
\lean{DeGiorgi.C_harnack}
\leanok
\uses{def:C_Moser, def:p0, def:CwH}
\begin{equation*}
  C_{\mathrm{Har}}(d) := 17\Bigl(\frac{d-2}{d-1}\,\log\bigl(C_{\mathrm{Moser}}(d)\,(d-1)^d\,4^d\bigr)
        + d + \log\bigl(C_{wH,0}(d)\,d^{\beta(d)}\bigr)\Bigr).
\end{equation*}
\end{definition}

\begin{lemma}
\label{lem:C_harnack_nonneg}
\lean{DeGiorgi.C_harnack_nonneg}
\leanok
\uses{def:Charnack}
For $d \geq 3$, $0 \leq C_{\mathrm{Har}}(d)$.
\end{lemma}

\begin{proof}
\leanok
\uses{def:C_Moser, def:p0, def:CwH, def:local-bases, lem:one_le_localMoserBase}

$(d-2)/(d-1) \geq 0$, both logs are nonnegative since
$M_{\mathrm{loc}}(d) \geq 1$ and $W_{\mathrm{loc}}(d) \geq 1$, and $d \geq 0$.
\end{proof}

\begin{theorem}
\label{thm:harnack}
\lean{DeGiorgi.harnack}
\leanok
\uses{def:ellipticityRatio, def:IsSubsolution, def:Charnack}
Let $d \geq 3$, $A$ normalized elliptic on $\Bz 1$, and $u > 0$ a (local
weak) solution of $-\dv(A\nabla u) = 0$ on $\Bz 1$. Then
\begin{equation*}
  \esssup_{\Bz{1/2}} u \;\leq\; \exp\bigl(C_{\mathrm{Har}}(d)\,\Lambda_A^{1/2}\bigr)
        \essinf_{\Bz{1/2}} u.
\end{equation*}
\end{theorem}

\begin{proof}
\leanok
\uses{def:MemW1p, def:MemW1pWitness, def:C_poinc_val, def:C_JN, lem:smoothTransition_nnnorm_deriv_bounded, thm:ae_on_set_of_ae_on_finite_cover, def:EllipticCoeff, def:rescaleToUnitBall, def:rescaleCoeffToUnitBall, thm:rescaleToUnitBall_isSubsolution, def:EllipticCoeff_restrict, def:moserChi, lem:moserDecayRatio_nonneg, def:C_Moser, lem:C_MoserAnchor_le_C_Moser, thm:linfty-Moser, def:C_weakHarnack0Forward, def:c_crossover_bmo_scale, lem:c-crossover-bounds, thm:crossover-exists, def:Ccrossover, def:p0, def:CwH, thm:wH-on-ball, def:local-bases, lem:one_le_localMoserBase, lem:meas-bd}

Define
\begin{equation*}
K(A) := \bigl(M_{\mathrm{loc}}(d)\Lambda_A^{d/2}\bigr)^{1/p_{\mathrm{Har}}(d)}
        \cdot W_{\mathrm{loc}}(d)^{\Lambda_A^{1/2}},
\end{equation*}
where $p_{\mathrm{Har}}(d) = (d-1)/(d-2) \cdot$ stuff (the auxiliary
exponent \leanname{harnack\_p}). The forward Moser estimate gives, on each
ball $\Ball{1/16}{c}$ centred at $c \in \Ball{1/2}{0}$,
\begin{equation*}
  \esssup_{\Ball{1/16}{c}} u
   \;\leq\; \bigl(M_{\mathrm{loc}}(d)\Lambda_A^{d/2}\bigr)^{1/p_{\mathrm{Har}}(d)}
        \cdot \Bigl(\fint_{\Ball{1/8}{c}} u^{p_{\mathrm{Har}}(d)}\Bigr)^{1/p_{\mathrm{Har}}(d)},
\end{equation*}
and Theorem~\ref{thm:wH-on-ball} (weak Harnack on $\Ball{1/8}{c}$) gives
\begin{equation*}
  \Bigl(\fint_{\Ball{1/8}{c}} u^{p_{\mathrm{Har}}(d)}\Bigr)^{1/p_{\mathrm{Har}}(d)}
        \;\leq\; W_{\mathrm{loc}}(d)^{\Lambda_A^{1/2}}\,
        \essinf_{\Ball{1/16}{c}} u.
\end{equation*}
Concatenating gives the local pointwise (a.e.) bound
$u(x) \leq K(A)\,\essinf_{\Ball{1/8}{c}} u$ for almost every $x \in \Ball{1/16}{c}$
(\leanname{ae\_le\_localHarnack\_on\_eighthBall}).

By compactness of $\overline{\Bz{1/2}}$, there is a finite cover of
$\overline{\Bz{1/2}}$ by balls $\Ball{1/16}{c_i}$ with $c_i = (15/16) x_i$,
$x_i \in \overline{\Bz{1/2}}$. Chaining at most $16$ such local Harnack
inequalities along a path connecting any two centres
(\leanname{quarterBallInf\_chain\_le}) yields, almost everywhere on
$\Bz{1/2}$,
\begin{equation*}
  u(x) \;\leq\; K(A)^{17}\,\essinf_{\Bz{1/2}} u.
\end{equation*}
Hence $\esssup_{\Bz{1/2}} u \leq K(A)^{17}\,\essinf_{\Bz{1/2}} u$.

Finally, the elementary identity
\begin{equation*}
  17\Bigl[\frac{d-2}{d-1}\log M_{\mathrm{loc}}(d) + \frac{d}{2}\log\Lambda_A
    + \log W_{\mathrm{loc}}(d) \Lambda_A^{1/2}\Bigr]
  \;\leq\; C_{\mathrm{Har}}(d)\,\Lambda_A^{1/2}
\end{equation*}
(using $\Lambda_A^{1/2} \geq 1$ and $\log\Lambda_A \leq 2\Lambda_A^{1/2}$,
established in \leanname{localHarnackConstant\_pow\_seventeen\_le\_exp})
gives $K(A)^{17} \leq \exp(C_{\mathrm{Har}}(d)\,\Lambda_A^{1/2})$, completing
the proof.
\end{proof}

\begin{corollary}
\label{cor:harnack_of_homogeneousWeakSolution}
\lean{DeGiorgi.harnack_of_homogeneousWeakSolution}
\leanok
\uses{def:ellipticityRatio, def:IsHomogeneousWeakSolution, def:Charnack}
The same conclusion holds for any positive homogeneous weak solution
$u$ on $\Bz 1$, since such a $u$ is in particular an
$\IsSolution$.
\end{corollary}

\begin{proof}
\leanok
\uses{lem:isHomogeneousWeakSolution_isSubsolution, thm:harnack}
See the Lean formalization.
\end{proof}

\subsection{Holder regularity}

\begin{definition}
\label{def:CholderMoser}
\lean{DeGiorgi.C_holder_Moser, DeGiorgi.moserDecayAlpha}
\leanok
\uses{def:ellipticityRatio, def:C_Moser, def:CwH, def:Charnack}
\begin{align*}
  C_{\mathrm{Hol}}(d) &:= 256\bigl(|C_{\mathrm{Har}}(d)| + M_{\mathrm{loc}}(d)
        + C_{\mathrm{Moser}}(d) 2^d + C_{wH,0}(d) d^d + 8\bigr), \\
  \alpha(A) &:= \exp\bigl(-(C_{\mathrm{Har}}(d)\,\Lambda_A^{1/2})\bigr).
\end{align*}
\end{definition}

\begin{lemma}
\label{lem:moserDecayAlpha_pos}
\lean{DeGiorgi.moserDecayAlpha_pos, DeGiorgi.moserDecayAlpha_le_one, DeGiorgi.moserLowerBound_le_decayAlpha}
\leanok
\uses{def:ellipticityRatio, def:CholderMoser}
$0 < \alpha(A) \leq 1$, and
$\exp(-(C_{\mathrm{Hol}}(d)\,\Lambda_A^{1/2})) \leq \alpha(A)$.
\end{lemma}

\begin{proof}
\leanok
\uses{def:C_poinc_val, def:C_JN, lem:smoothTransition_nnnorm_deriv_bounded, def:moserChi, lem:moserDecayRatio_nonneg, def:C_Moser, lem:C_MoserAnchor_le_C_Moser, def:C_weakHarnack0Forward, def:c_crossover_bmo_scale, lem:c-crossover-bounds, thm:crossover-exists, def:Ccrossover, def:CwH, def:local-bases, def:Charnack, lem:C_harnack_nonneg}

$\alpha(A) > 0$ as the exponential of a real, $\alpha(A) \leq 1$ since the
exponent is $\leq 0$ (because $C_{\mathrm{Har}} \geq 0$), and the lower
bound is the monotonicity of $\exp$ together with $C_{\mathrm{Har}}(d) \leq
C_{\mathrm{Hol}}(d)$ (\leanname{C\_harnack\_le\_C\_holder\_Moser}).
\end{proof}

\begin{lemma}
\label{lem:one_sub_moserDecayAlpha_le_two_rpow_neg}
\lean{DeGiorgi.one_sub_moserDecayAlpha_le_two_rpow_neg}
\leanok
\uses{def:ellipticityRatio, def:CholderMoser}
$1 - \alpha(A) \leq (1/2)^{\alpha(A)}$.
\end{lemma}

\begin{proof}
\leanok
\uses{lem:moserDecayAlpha_pos}

Set $\alpha = \alpha(A) \in (0,1]$. We have $1 - \alpha \leq e^{-\alpha}$
(by $1+x \leq e^x$ at $x = -\alpha$), and $-\alpha \leq -\alpha\log 2$ since
$\log 2 \leq 1$. Hence $e^{-\alpha} \leq e^{-\alpha\log 2} = (1/2)^\alpha$.
\end{proof}

\begin{theorem}
\label{thm:holder-Moser}
\lean{DeGiorgi.holder_Moser}
\leanok
\uses{def:ellipticityRatio, def:IsSubsolution, def:CholderMoser}
Let $d \geq 3$, $A$ normalized elliptic on $\Bz 1$, $p_0 > 1$, and $u$ a
solution of $-\dv(A\nabla u) = 0$ with $|u|^{p_0} \in \Lp{1}(\Bz 1)$.
Then there exists a representative $v$ of $u$ on $\Bz 1$ (i.e.\ $v = u$
almost everywhere on $\Bz 1$) and an exponent
$\alpha = \alpha(A) > 0$ with
$\exp(-(C_{\mathrm{Hol}}(d)\,\Lambda_A^{1/2})) \leq \alpha$ such that for
every $x, y \in \Bz{1/2}$,
\begin{equation*}
  |v(x) - v(y)|
   \;\leq\; C_{\mathrm{Hol}}(d)\,\Lambda_A^{d/(2p_0)}
        \Bigl(\frac{p_0}{p_0 - 1}\Bigr)^{d/p_0}
        \Bigl(\int_{\Bz 1} |u|^{p_0}\,dz\Bigr)^{1/p_0}
        \,\|x - y\|^{\alpha}.
\end{equation*}
\end{theorem}

\begin{proof}
\leanok
\uses{def:MemW1p, def:MemW1pWitness, def:C_poinc_val, def:C_JN, def:dyadicBallAverage, lem:smoothTransition_nnnorm_deriv_bounded, def:rescaleToUnitBall, def:rescaleCoeffToUnitBall, thm:rescaleToUnitBall_isSubsolution, def:EllipticCoeff_restrict, lem:essSup_rescale_halfBall, def:moserChi, lem:moserDecayRatio_nonneg, def:C_Moser, lem:C_MoserAnchor_le_C_Moser, thm:linfty-Moser, def:C_weakHarnack0Forward, def:c_crossover_bmo_scale, lem:c-crossover-bounds, thm:crossover-exists, def:Ccrossover, def:CwH, def:local-bases, lem:meas-bd, def:Charnack, lem:C_harnack_nonneg, thm:harnack, lem:moserDecayAlpha_pos, lem:one_sub_moserDecayAlpha_le_two_rpow_neg}

Define $\alpha := \alpha(A)$ and let $v := \leanname{moserRepresentative}\,u$
be the dyadic-average representative supplied by
\leanname{Holder/Representative.lean}. By
\leanname{moserRepresentative\_ae\_eq}, $v = u$ almost everywhere on $\Bz 1$.

Fix $x, y \in \Bz{1/2}$. If $x = y$ the bound is trivial. Otherwise,
distinguish two cases.

\emph{Large distance.} If $\|x - y\| \geq 1/8$, by
\leanname{moserRepresentative\_large\_distance\_bound} there is a direct
control by the $L^{p_0}$ norm of $u$:
\begin{equation*}
  |v(x) - v(y)| \;\leq\; C_{\mathrm{Hol}}(d)\,\Lambda_A^{d/(2p_0)}
    \bigl(p_0/(p_0-1)\bigr)^{d/p_0}\,\|u\|_{L^{p_0}(\Bz 1)}.
\end{equation*}
Since $\|x-y\|^\alpha \geq (1/8)^\alpha \geq 1/8$ and absorbing constants
into $C_{\mathrm{Hol}}$ via the factor $256$ in
Definition~\ref{def:CholderMoser} gives the claim.

\emph{Small distance.} If $\|x - y\| < 1/8$, choose $n$ with
$2^{-(n+1)} \leq \|x - y\| < 2^{-n}$ \leanname{(exists\_moserDyadicRadius\_near)},
and apply the dyadic oscillation decay
\leanname{moser\_oscillation\_decay\_on\_ball}, which iterates the Harnack
inequality on each ball $\Ball{2^{-k}}{c}$ centred near $x$:
\begin{equation*}
  \osc_{\Ball{2^{-k-1}}{c}} u
    \;\leq\; (1 - \alpha)\,\osc_{\Ball{2^{-k}}{c}} u.
\end{equation*}
Iterating gives geometric decay
$\osc_{\Ball{2^{-n}}{c}} u \leq (1-\alpha)^n\,\osc_{\Ball 1 0} u$.
By the previous lemma $1 - \alpha \leq (1/2)^\alpha$, hence
$(1-\alpha)^n \leq 2^{-\alpha n} \leq C\,\|x - y\|^\alpha$.

The starting oscillation $\osc_{\Bz 1} u$ is controlled by twice the local
$L^\infty$ bound, which by the forward Moser estimate
\leanname{linfty\_subsolution\_Moser\_on\_ball} is bounded by
$M_{\mathrm{loc}}(d)\Lambda_A^{d/(2p_0)} (p_0/(p_0-1))^{d/p_0}
        \|u\|_{L^{p_0}(\Bz 1)}$.
This yields the desired pointwise Holder bound on the representative $v$.
\end{proof}

\begin{corollary}
\label{cor:holder_Moser_of_homogeneousWeakSolution}
\lean{DeGiorgi.holder_Moser_of_homogeneousWeakSolution}
\leanok
\uses{def:ellipticityRatio, def:IsHomogeneousWeakSolution, def:CholderMoser}
The same conclusion holds verbatim for homogeneous weak solutions
($\IsHomogeneousWeakSolution$), since these satisfy
$\IsSolution$.
\end{corollary}

\begin{proof}
\leanok
\uses{lem:isHomogeneousWeakSolution_isSubsolution, thm:holder-Moser}
See the Lean formalization.
\end{proof}

\begin{remark}
The dependence on the ellipticity bound $\Lambda_A$ enters Holder regularity
in two ways: through the multiplicative prefactor $\Lambda_A^{d/(2p_0)}$
inherited from the forward Moser bound, and through the Holder exponent
$\alpha(A) = \exp(-C_{\mathrm{Har}}(d)\,\Lambda_A^{1/2})$ inherited from the
Harnack constant. As $\Lambda_A \to \infty$, $\alpha(A) \to 0$ exponentially,
quantifying the well-known degeneration of De Giorgi--Nash--Moser regularity
for highly anisotropic elliptic operators.
\end{remark}
